
\subsection{{Constraints}}

\purple{
	TRATAMOS ISSO NO CAPITULO ANTERIOR COMO UMA ANALISE MAIS GERAL DAS RESTRICOES, PRECISAMOS 
	UNIFORMIZAR OS TERMOS AQUI? 

	As restricoes do controlador devem descrever os limites de operacao da dinamica do modelo e ainda encontrar uma solu��o que leve o sistema
	para mais proximo da referencia.

	\red{ESCREVER EQUACOES DE RESTRICAO}

	\red{descrever que a trajetoria deve ser admissivel e forma com que o sistema calcula o controle tambem deve ser admissivel}

	\red{fazer um rapido comentario sobre a forma que essa restricao � trabalhada no capitulo 2}
}

\subsection{{Feasibility region}}

\red{nao aplicavel ????}

\red{como calcular ???}
	restricoes do problem de prog quad

	sobre o controle

	variacoes do instante de chaveadmento

	mas envolvem o estado inicial do ciclo

	sea dependencia das restricoes com respeito ao estao inicial for linear

	entao podemos encarar o x0 como uma variavel junto com as variaveis de controle

	com isso as restricoes descrevem um polietro aumenta do de controle e estado inicial

	geometria computacional mpttoolbox

	projeta a geometria no subespaco

\red{pode haver melhoria com o ponto final livre}

\red{interessante para avaliar o desempenho do controlador na regiao de factibilidade}

\textbf{APPLICATION EXAMPLES}

\begin{itemize}
	\item Integrador duplo( matheus )
	\item Patino 1
	\item Patino 2
\end{itemize}